%%%%%%%%%%%%%%%%%%%%%%%%%%%%%%%%%%%%%%%%%%%%%%%%%%%%%%%%%%%%%%%%%%%%%%%%%%%%%%

%\section{Polynomial bound for the number of mixing operations}
%\label{sec: polynomial proof}
%%%%%%%%%%%%%%%%%%%%%%%%%%%%%%%%%%%%%%%%%%%%%%%%%%%%%%%%%%%%%%%%%%%%%%%%%%%%%%%

%\section{Polynomial bound for the number of mixing operations}
%\label{sec: polynomial proof}
%%%%%%%%%%%%%%%%%%%%%%%%%%%%%%%%%%%%%%%%%%%%%%%%%%%%%%%%%%%%%%%%%%%%%%%%%%%%%%%

%\section{Polynomial bound for the number of mixing operations}
%\label{sec: polynomial proof}
%%%%%%%%%%%%%%%%%%%%%%%%%%%%%%%%%%%%%%%%%%%%%%%%%%%%%%%%%%%%%%%%%%%%%%%%%%%%%%%

%\section{Polynomial bound for the number of mixing operations}
%\label{sec: polynomial proof}
%\input{6_polynomial_proof.tex}

Let $C\ingroundset\integers$ with $\average(C)\in\integers$ and 
$\nofdroplets{C}=n \ge 4$ be a configuration that satisfies Condition~$\MixCond$.
The existence of a perfect-mixing graph for $C$ was established in 
Section~\ref{sec: sufficiency of condition mc}. This graph, however, might be very
large -- it can be shown that if arbitrary droplets are mixed at each step then
it might take an exponential number of steps for the process to converge.
In this section we prove Theorem~\ref{thm: miscibility characterization}(b),
namely that $C$ can be perfectly mixed
with precision at most $1$ and in a polynomial number of mixing operations. 
The essence of the proof is
to show that in the construction in  Section~\ref{sec: sufficiency of condition mc}
it is possible to choose a mixing operation at each step so that the overall
number of steps will be polynomial in the input size.
We assume here that the reader is familiar with the results from  Section~\ref{sec: sufficiency of condition mc};
in fact, some of the lemmas or observations from that section may be used here occasionally
without an explicit reference.

It is sufficient to show that configuration $E$, constructed from $C$ 
in Section~\ref{subsec: proof for n greater-equal than 7}, 
is perfectly mixable with precision $0$ in a polynomial number of mixing operations.
(As described in Section~\ref{subsec: proof for n greater-equal than 7},
constructing such $E$ from $C$ requires only two mixing operations. Recall that this
construction also involves a linear mapping, but this mapping does not change the
constructed mixing graph.)
For this reason we will assume in this section that $E$ is the initial configuration.

Recall that $\potential(E)=\sum_{e\in E}(e-\hatmu)^2$, where
$\hatmu =  \average(E)\in\integers$. We use $\potential(E)$ to measure the progress of
the mixing process, ultimately showing that $\potential(E)$
can be decreased down to $0$ after a number of steps that is polynomial in  
the initial value of $\log\potential(E)$, and thus also in $s(E)$, the size of $E$.
(What we actually show is that after the mixing process we achieve $\potential(E)\leq 1$. 
Since $E\cup \braced{\hatmu} \subseteq \integers$, this implies that in fact $\potential(E)=0$.)

The general idea is to always mix two concentrations whose difference 
is large enough, so that after a polynomial number of mixing operations,
$\potential(E)$ decreases at least by a factor of $\half$. This is sufficient to 
establish a polynomial bound for the whole process.  Of course, the  pair of
concentrations that we mix must either preserve the corresponding invariant
or produce a near-final configuration.
When $n$ is a power of $2$, it is relatively
simple to identify good pairs to mix (see Section~\ref{sec: polynomial - near final}), 
because then we only need to ensure that the 
concentrations mixed at each step have the same parity. This also easily
extends to near-final configurations.
For arbitrary configurations $E$ (that satisfy the appropriate invariant), however,
identifying such good pairs to mix is more challenging (see 
Sections~\ref{sec: an exponential bound} and~\ref{sec: polynomial - arbitrary n}).

In the following sub-sections, we will present some auxiliary mixing sequences 
that will be later combined to construct a polynomial-length mixing sequence for $E$.
To streamline the arguments, we will focus on estimating the length of these sequences;
that these mixing sequences can actually be computed in polynomial time will be implicit 
in their construction. We will return to the time complexity analysis
in Section~\ref{sec: polynomial running time}.





Let $C\ingroundset\integers$ with $\average(C)\in\integers$ and 
$\nofdroplets{C}=n \ge 4$ be a configuration that satisfies Condition~$\MixCond$.
The existence of a perfect-mixing graph for $C$ was established in 
Section~\ref{sec: sufficiency of condition mc}. This graph, however, might be very
large -- it can be shown that if arbitrary droplets are mixed at each step then
it might take an exponential number of steps for the process to converge.
In this section we prove Theorem~\ref{thm: miscibility characterization}(b),
namely that $C$ can be perfectly mixed
with precision at most $1$ and in a polynomial number of mixing operations. 
The essence of the proof is
to show that in the construction in  Section~\ref{sec: sufficiency of condition mc}
it is possible to choose a mixing operation at each step so that the overall
number of steps will be polynomial in the input size.
We assume here that the reader is familiar with the results from  Section~\ref{sec: sufficiency of condition mc};
in fact, some of the lemmas or observations from that section may be used here occasionally
without an explicit reference.

It is sufficient to show that configuration $E$, constructed from $C$ 
in Section~\ref{subsec: proof for n greater-equal than 7}, 
is perfectly mixable with precision $0$ in a polynomial number of mixing operations.
(As described in Section~\ref{subsec: proof for n greater-equal than 7},
constructing such $E$ from $C$ requires only two mixing operations. Recall that this
construction also involves a linear mapping, but this mapping does not change the
constructed mixing graph.)
For this reason we will assume in this section that $E$ is the initial configuration.

Recall that $\potential(E)=\sum_{e\in E}(e-\hatmu)^2$, where
$\hatmu =  \average(E)\in\integers$. We use $\potential(E)$ to measure the progress of
the mixing process, ultimately showing that $\potential(E)$
can be decreased down to $0$ after a number of steps that is polynomial in  
the initial value of $\log\potential(E)$, and thus also in $s(E)$, the size of $E$.
(What we actually show is that after the mixing process we achieve $\potential(E)\leq 1$. 
Since $E\cup \braced{\hatmu} \subseteq \integers$, this implies that in fact $\potential(E)=0$.)

The general idea is to always mix two concentrations whose difference 
is large enough, so that after a polynomial number of mixing operations,
$\potential(E)$ decreases at least by a factor of $\half$. This is sufficient to 
establish a polynomial bound for the whole process.  Of course, the  pair of
concentrations that we mix must either preserve the corresponding invariant
or produce a near-final configuration.
When $n$ is a power of $2$, it is relatively
simple to identify good pairs to mix (see Section~\ref{sec: polynomial - near final}), 
because then we only need to ensure that the 
concentrations mixed at each step have the same parity. This also easily
extends to near-final configurations.
For arbitrary configurations $E$ (that satisfy the appropriate invariant), however,
identifying such good pairs to mix is more challenging (see 
Sections~\ref{sec: an exponential bound} and~\ref{sec: polynomial - arbitrary n}).

In the following sub-sections, we will present some auxiliary mixing sequences 
that will be later combined to construct a polynomial-length mixing sequence for $E$.
To streamline the arguments, we will focus on estimating the length of these sequences;
that these mixing sequences can actually be computed in polynomial time will be implicit 
in their construction. We will return to the time complexity analysis
in Section~\ref{sec: polynomial running time}.





Let $C\ingroundset\integers$ with $\average(C)\in\integers$ and 
$\nofdroplets{C}=n \ge 4$ be a configuration that satisfies Condition~$\MixCond$.
The existence of a perfect-mixing graph for $C$ was established in 
Section~\ref{sec: sufficiency of condition mc}. This graph, however, might be very
large -- it can be shown that if arbitrary droplets are mixed at each step then
it might take an exponential number of steps for the process to converge.
In this section we prove Theorem~\ref{thm: miscibility characterization}(b),
namely that $C$ can be perfectly mixed
with precision at most $1$ and in a polynomial number of mixing operations. 
The essence of the proof is
to show that in the construction in  Section~\ref{sec: sufficiency of condition mc}
it is possible to choose a mixing operation at each step so that the overall
number of steps will be polynomial in the input size.
We assume here that the reader is familiar with the results from  Section~\ref{sec: sufficiency of condition mc};
in fact, some of the lemmas or observations from that section may be used here occasionally
without an explicit reference.

It is sufficient to show that configuration $E$, constructed from $C$ 
in Section~\ref{subsec: proof for n greater-equal than 7}, 
is perfectly mixable with precision $0$ in a polynomial number of mixing operations.
(As described in Section~\ref{subsec: proof for n greater-equal than 7},
constructing such $E$ from $C$ requires only two mixing operations. Recall that this
construction also involves a linear mapping, but this mapping does not change the
constructed mixing graph.)
For this reason we will assume in this section that $E$ is the initial configuration.

Recall that $\potential(E)=\sum_{e\in E}(e-\hatmu)^2$, where
$\hatmu =  \average(E)\in\integers$. We use $\potential(E)$ to measure the progress of
the mixing process, ultimately showing that $\potential(E)$
can be decreased down to $0$ after a number of steps that is polynomial in  
the initial value of $\log\potential(E)$, and thus also in $s(E)$, the size of $E$.
(What we actually show is that after the mixing process we achieve $\potential(E)\leq 1$. 
Since $E\cup \braced{\hatmu} \subseteq \integers$, this implies that in fact $\potential(E)=0$.)

The general idea is to always mix two concentrations whose difference 
is large enough, so that after a polynomial number of mixing operations,
$\potential(E)$ decreases at least by a factor of $\half$. This is sufficient to 
establish a polynomial bound for the whole process.  Of course, the  pair of
concentrations that we mix must either preserve the corresponding invariant
or produce a near-final configuration.
When $n$ is a power of $2$, it is relatively
simple to identify good pairs to mix (see Section~\ref{sec: polynomial - near final}), 
because then we only need to ensure that the 
concentrations mixed at each step have the same parity. This also easily
extends to near-final configurations.
For arbitrary configurations $E$ (that satisfy the appropriate invariant), however,
identifying such good pairs to mix is more challenging (see 
Sections~\ref{sec: an exponential bound} and~\ref{sec: polynomial - arbitrary n}).

In the following sub-sections, we will present some auxiliary mixing sequences 
that will be later combined to construct a polynomial-length mixing sequence for $E$.
To streamline the arguments, we will focus on estimating the length of these sequences;
that these mixing sequences can actually be computed in polynomial time will be implicit 
in their construction. We will return to the time complexity analysis
in Section~\ref{sec: polynomial running time}.





Let $C\ingroundset\integers$ with $\average(C)\in\integers$ and 
$\nofdroplets{C}=n \ge 4$ be a configuration that satisfies Condition~$\MixCond$.
The existence of a perfect-mixing graph for $C$ was established in 
Section~\ref{sec: sufficiency of condition mc}. This graph, however, might be very
large -- it can be shown that if arbitrary droplets are mixed at each step then
it might take an exponential number of steps for the process to converge.
In this section we prove Theorem~\ref{thm: miscibility characterization}(b),
namely that $C$ can be perfectly mixed
with precision at most $1$ and in a polynomial number of mixing operations. 
The essence of the proof is
to show that in the construction in  Section~\ref{sec: sufficiency of condition mc}
it is possible to choose a mixing operation at each step so that the overall
number of steps will be polynomial in the input size.
We assume here that the reader is familiar with the results from  Section~\ref{sec: sufficiency of condition mc};
in fact, some of the lemmas or observations from that section may be used here occasionally
without an explicit reference.

It is sufficient to show that configuration $E$, constructed from $C$ 
in Section~\ref{subsec: proof for n greater-equal than 7}, 
is perfectly mixable with precision $0$ in a polynomial number of mixing operations.
(As described in Section~\ref{subsec: proof for n greater-equal than 7},
constructing such $E$ from $C$ requires only two mixing operations. Recall that this
construction also involves a linear mapping, but this mapping does not change the
constructed mixing graph.)
For this reason we will assume in this section that $E$ is the initial configuration.

Recall that $\potential(E)=\sum_{e\in E}(e-\hatmu)^2$, where
$\hatmu =  \average(E)\in\integers$. We use $\potential(E)$ to measure the progress of
the mixing process, ultimately showing that $\potential(E)$
can be decreased down to $0$ after a number of steps that is polynomial in  
the initial value of $\log\potential(E)$, and thus also in $s(E)$, the size of $E$.
(What we actually show is that after the mixing process we achieve $\potential(E)\leq 1$. 
Since $E\cup \braced{\hatmu} \subseteq \integers$, this implies that in fact $\potential(E)=0$.)

The general idea is to always mix two concentrations whose difference 
is large enough, so that after a polynomial number of mixing operations,
$\potential(E)$ decreases at least by a factor of $\half$. This is sufficient to 
establish a polynomial bound for the whole process.  Of course, the  pair of
concentrations that we mix must either preserve the corresponding invariant
or produce a near-final configuration.
When $n$ is a power of $2$, it is relatively
simple to identify good pairs to mix (see Section~\ref{sec: polynomial - near final}), 
because then we only need to ensure that the 
concentrations mixed at each step have the same parity. This also easily
extends to near-final configurations.
For arbitrary configurations $E$ (that satisfy the appropriate invariant), however,
identifying such good pairs to mix is more challenging (see 
Sections~\ref{sec: an exponential bound} and~\ref{sec: polynomial - arbitrary n}).

In the following sub-sections, we will present some auxiliary mixing sequences 
that will be later combined to construct a polynomial-length mixing sequence for $E$.
To streamline the arguments, we will focus on estimating the length of these sequences;
that these mixing sequences can actually be computed in polynomial time will be implicit 
in their construction. We will return to the time complexity analysis
in Section~\ref{sec: polynomial running time}.



